%% Complete proofs for the derived negative translation calculus.
\section{Negative Free Logic: Translation and Guarded Rules}
\label{sec:annex-p08-negative-free-logic-calculus}

\begin{annexitem}{annex:p8.negative.model-bridge}{thm:p08-negative-model-bridge}{Value induction, truth induction, encoding and decoding}
\begin{proof}
First work in a fixed negative model with an outer assignment.  Encode its
outer domain as the classical domain, interpret the new predicate by \(D\),
and keep the total functions and strict ordinary predicate extensions.
The target theory \(\Theta\) is true because \(D\ne\varnothing\).

We prove the value and truth assertions simultaneously by the size of a
finite term/formula syntax tree.  A variable has its assigned outer value;
\(z=x\) holds exactly for that value.  The constant case is the same using
its interpretation.  For \(f(t_1,\ldots,t_n)\), the induction hypothesis
identifies each existentially supplied \(u_i\) with the defined value of
\(t_i\).  Such a tuple exists exactly when every argument is defined.
When it exists its entries are unique and the total function has exactly
one value; the last target equality fixes \(z\) to that value.  If some
argument is undefined, no tuple satisfies the conjunction, in agreement
with the strict lifting.  For arity zero the function is a constant and
the displayed empty-tuple convention gives the same conclusion.  This
also proves that all terms of the ordinary language are defined by
induction, without asserting that their values are in \(D\).

For a description \(\iota xB\), the formula induction hypothesis on its
proper constituent \(B\) identifies its satisfying existing objects with
the objects satisfying \(E(u)\land\delta B[u/x]\).  Fresh auxiliary
variables make target substitution equivalent to changing the one
assignment coordinate.  The value formula states that \(z\) is one of
these existing objects and every such object equals \(z\).  It therefore
has exactly one value when the satisfying set has exactly one member,
and no value when that set is empty or has at least two members.  This
is exactly the primitive description evaluation.  These cases establish
the uniqueness assertion for partial terms as well.

At a positive-arity predicate atom, the term assertions identify the
existentially quantified argument tuple with its unique defined tuple.
The \(E\)-conjuncts require every complete argument value to lie in
\(D\); the remaining predicate conjunct makes precisely the source
extension test.  Undefined and defined-but-nonexisting arguments both
make the target atom false, as required.  A nullary predicate has no
argument to test, and its target and source Boolean values coincide.
The existence atom has its one value guard and membership guard, so it
is true exactly for a defined existing value.  In the identity case the
two value guards, two membership guards and target equality require
precisely the same existing object.  Thus target equality is never used
to make the source identity of a nonexisting term true.

Falsum is false in both languages.  Negation exchanges true and false,
and conjunction, disjunction and implication apply their same two-valued
truth functions to the constituent values; the induction hypothesis
proves each Boolean clause.  For a universal formula the target
quantifier tests every object of \(U\), but an object outside \(D\)
makes its implication true vacuously.  At objects of \(D\), the formula
induction hypothesis says that the consequent holds exactly when the
source body holds after the same assignment change.  Hence the target
universal is true exactly when every source inner-domain instance is
true.  The existential clause similarly requires a target witness in
\(D\), at which the body agrees; this is exactly a source existential
witness.  Bound alpha-renaming keeps all these assignment changes
independent of the freshly quantified value variables.  All recursive
uses are on proper constituents, completing the simultaneous induction.

For the converse model direction, start with a classical model of
\(\Theta\).  Its domain \(U\) is nonempty by ordinary first-order
semantics, and \(D=E\)'s extension is nonempty by \(\Theta\).
Keep its constants and total functions, and replace every predicate
extension by its intersection with \(D^n\), leaving nullary values
unchanged.  This produces an actual negative model.  The same induction
works: at a predicate atom the target guards ensure that only tuples
in \(D^n\) are tested, exactly those retained by restriction.  It is
irrelevant that the target predicate might originally have held on
additional outer tuples.  Thus this construction is a truth-preserving
decoding, not an assertion that encoding and decoding are inverse on
irrelevant predicate data.

Finally suppose negative consequence holds.  Any target model and
assignment satisfying \(\Theta\), \(\mathcal E_X\) and
\(\delta[\Gamma]\) decodes to a negative model with the same assignment
satisfying the parameter restrictions and all source premises.  Its
source conclusion is true, hence its target conclusion is true by the
proved induction.  This proves the target consequence.  Conversely,
encode any negative model and permitted assignment satisfying
\(\Gamma\).  The target premises hold, target consequence gives
\(\delta A\), and the induction gives source truth of \(A\).
This proves both directions for arbitrary premise and parameter sets.
\end{proof}
\end{annexitem}

\begin{annexitem}{annex:p8.negative.substitution}{thm:p08-negative-substitution}{Defined-term substitution including nested description bodies}
\begin{proof}
Let the defined value of \(t\) under \(g\) be \(a\).  Rename every binder
in the substituted expression to avoid the finitely many free variables of
\(t\), as well as the substituted variable wherever necessary.  Prove
the two assertions simultaneously by syntax-tree induction.

For a variable equal to \(x\), substitution inserts \(t\), whose value
is \(a\), and changing the assignment gives the same \(a\).  For a
different variable the assignment is unchanged; constants are unchanged.
At a function application each argument has equal partial values on the
two sides by induction.  If any is undefined, both applications are
undefined.  Otherwise their equal value tuples have the same image
under the same total function.  For a description with fresh binder
\(y\), \(t\)'s value is unchanged when \(g\) is modified at \(y\).
For every \(d\in D\), induction on its body identifies truth under
\(g[y\mapsto d]\) after substitution with truth under
\(g[x\mapsto a,y\mapsto d]\).  The two modifications commute because
the binder has been freshened.  The descriptions therefore have the
same set of satisfying existing objects; its singleton, empty and
multiple-member cases give equal partial values.

Predicate, identity and existence atoms test equal partial argument
values on the two sides, so their Boolean truth values agree, including
their false undefined/nonexisting cases.  Nullary atoms are unchanged.
The same classical truth functions preserve agreement under each
Boolean connective and falsum.  For a fresh universal binder \(y\),
the defined value of \(t\) remains \(a\) for each inner-domain change
at \(y\); the formula induction hypothesis identifies every quantified
instance on the two sides, and hence the universal truth values.
The existential case uses the identical instance agreement and
identifies the existence of a witness.  A binder originally named \(x\)
has been alpha-renamed, or else masks substitution at exactly its bound
occurrences, giving the same result.  These cases exhaust the mutually
generated grammar.  None used \(a\in D\): definedness alone suffices
for substitution; inner-domain membership is the further quantifier
guard used in the next proof.
\end{proof}
\end{annexitem}

\begin{annexitem}{annex:p8.negative.adequacy}{thm:p08-negative-adequacy}{Strong completeness via an actual model decoding}
\begin{proof}
If a certificate exists, ordinary first-order natural-deduction soundness,
Theorem~\ref{thm:fol-soundness}, makes its target conclusion a
consequence of the target theory, parameter guards and translated premises.
Theorem~\ref{thm:p08-negative-model-bridge} then gives the exact negative
consequence.  This establishes soundness for every allowed signature.

Conversely, if \(\Gamma\models_{-,X}A\), the model bridge gives
\(\Theta\cup\mathcal E_X\cup\delta[\Gamma]\models\delta A\)
in ordinary classical first-order logic.  For countable signatures the
target signature is still countable: it adds only one predicate and an
infinite variable supply, and translation introduces no other function
or predicate symbols.  Corollary~\ref{cor:fol-completeness-countable}
therefore supplies a finite proof on the registered ND-FOL sheet from
those target assumptions, without Choice.  For arbitrary signatures
Theorem~\ref{thm:fol-completeness} supplies the same proof under its
stated ambient Choice assumption.  In either case the proof is precisely
the certificate required by Definition~\ref{def:p08-negative-calculus}.
No step assumes that the native shortcuts already form a complete system.

A finite certificate has finitely many open assumption occurrences.
Collect their source-premise preimages into \(\Gamma_0\) and their
parameter names into \(X_0\), both finite.  Its other open assumptions
are drawn from the singleton \(\Theta\).  It is then a certificate
relative to \(\Gamma_0,X_0\); soundness gives the corresponding finite
negative consequence.  Conversely such a finite consequence remains
true when further premises or existing-parameter requirements are added.
This proves the finite-support and compactness assertions.  For sentences,
the truth induction is assignment-independent, so inner versus outer
free-variable conventions change neither side of the equivalence.
\end{proof}
\end{annexitem}

\begin{annexitem}{annex:p8.negative.rules}{thm:p08-negative-native-rules}{Every registered native interface and its freshness conditions}
\begin{proof}
The adequacy theorem has already been proved using only certificates,
the value/truth construction and ordinary first-order metatheory.
For each native row it therefore suffices to prove consequence
preservation, after which adequacy supplies the conclusion certificate.
The defining certificate row itself is sound by the preceding proof.

An assumption holds whenever all its displayed premises hold.
Implication elimination and conjunction introduction/eliminations follow
from their classical truth functions at one permitted assignment.
Either disjunction introduction makes its disjunction true when the given
disjunct is true. For disjunction elimination, a true major disjunction
has a true disjunct; the corresponding branch consequence then gives the
common conclusion while the remaining hypotheses are retained.
For implication introduction, if \(A\) is false the conclusion is true;
if \(A\) is true the premise consequence gives \(B\), making the
conclusion true.  Negation introduction similarly excludes a true \(A\)
when the assumed premise would imply falsum.  Negation elimination
cannot have both premises true.  Falsum elimination has no model
satisfying its premise.  For classical reductio, if its conclusion
were false then its negation would be true, and the premise consequence
would give falsum, a contradiction.  All these arguments leave \(X\)
and the other hypotheses unchanged.  Existing-parameter introduction
is sound by the defining requirement on \(g(x)\) for \(x\in X\).

For guarded universal elimination, truth of \(E!(t)\) supplies a defined
value \(a\in D\).  The universal premise gives the body under
\(g[x\mapsto a]\).  The substitution lemma gives the displayed
instance under \(g\).  Guarded existential introduction reverses this
argument: its instance and the same existence guard give a body true
at an inner witness, hence the existential conclusion.

For guarded universal introduction, fix any permitted assignment
satisfying \(\Gamma\), and any \(a\in D\).  Change the fresh
eigenvariable \(y\) to \(a\).  Since \(y\notin X\), this leaves all
parameter requirements intact.  Since \(y\) is absent from the other
hypotheses, they remain true; the temporary existence premise is
also true.  The premise consequence yields \(A[y/x]\).
The substitution lemma for the variable \(y\) identifies this with
the body at \(x=a\).  Freshness from the quantified conclusion makes
the truth of that conclusion independent of the changed \(y\).
As \(a\) was arbitrary in \(D\), the universal conclusion holds.

For guarded existential elimination, fix an assignment satisfying
\(\Gamma\cup\Delta\).  The major consequence gives an inner-domain
witness \(a\) for the existential formula.  Change fresh \(y\) to
\(a\).  Neither the parameter requirements nor any other hypothesis
changes.  Its existence and substituted-body assumptions now hold by
the witness and substitution lemma.  The minor consequence gives
\(B\), and freshness of \(y\) from \(B\) gives the same truth under
the original assignment.  This proves the conclusion while explaining
both temporary hypotheses and every displayed freshness condition.

Guarded strict reflexivity holds because the existence premise supplies
one existing value, equal to itself.  A strict identity premise supplies
defined existing values on both sides, so both identity-to-existence
rows hold.  In replacement those values are the same object \(a\);
the substitution lemma identifies both displayed formula instances
with truth of the same formula at \(x=a\), proving preservation even
through quantifiers and description bodies.  Truth of an ordinary
positive-arity atom requires every complete argument value in \(D\),
giving the atomic-existence row.  It makes no assertion about a proper
subterm, whose total ordinary function may map an outer value into \(D\).

For the nonempty-domain row choose one \(a\in D\), possible because
the model convention excludes an empty inner domain.  Change fresh
\(y\) to \(a\).  Its absence from \(X,\Gamma,A\) preserves every
parameter condition, every original hypothesis and the conclusion.
The temporary existence hypothesis holds, so the premise consequence
gives \(A\) at the modified and then the original assignment.
Thus this rule depends exactly on nonemptiness and freshness, not on
existence of any chosen ordinary constant.  Every displayed native
row has now been checked, and independent certificate completeness
establishes its admissibility.
\end{proof}
\end{annexitem}

\begin{annexitem}{annex:p8.negative.equality-existence}{thm:p08-negative-equality-existence}{Both guarded equivalences and the nonempty existential theorem}
\begin{proof}
From \(E!(t)\), guarded strict reflexivity gives \(t=t\).
From \(t=t\), identity-to-existence gives \(E!(t)\).
Using \(t=t\) with the supplied existence guard in existential
introduction gives \(\exists x(x=t)\), for fresh \(x\).
Conversely eliminate \(\exists x(x=t)\) using a fresh witness
\(y\), temporarily assuming \(E!(y)\) and \(y=t\).
Identity-to-existence gives \(E!(t)\); guarded existential
elimination removes the temporary hypotheses.  These are all
derived rows whose admissibility was proved in the preceding item,
and give each direction as a finite negative derivability.

For \(\exists xE!(x)\), temporarily assume \(E!(y)\) for a
fresh variable.  Guarded existential introduction, with the same
formula both as instance and guard, gives the existential formula.
Apply the nonempty-domain row to remove \(E!(y)\).
The nonexisting-constant countermodel in the next item shows why
this does not supply unconditional \(t=t\) for a prescribed term.
\end{proof}
\end{annexitem}

\begin{annexitem}{annex:p8.negative.boundaries}{thm:p08-negative-boundaries}{Nonexisting values, proper subterms, parameters, and failed descriptions}
\begin{proof}
In the two-object model the only quantified object is \(0\).
It exists, so \(\forall xE!(x)\) is true and
\(\exists x\neg E!(x)\) false.  The constant denotes \(1\),
outside \(D\), so its existence and strict self-identity are
false and its negated existence true.  Universal instantiation
without a guard would take the true universal to false \(E!(c)\);
existential generalization without a guard would take true
\(\neg E!(c)\) to false \(\exists x\neg E!(x)\).
Unconditional reflexivity would assert false \(c=c\).

For the function example define its other value arbitrarily in
\(U\), for instance \(f(0)=0\), making a total ordinary function.
Its value at \(c\) is \(0\), an existing object in the predicate
extension, so \(P(f(c))\) is true.  Nevertheless the proper
subterm \(c\) still denotes the nonexisting object \(1\).
This satisfies every declared atomic/function clause and refutes
the proposed extraction of existence for all proper subterms.
Similarly choosing \(f(0)=1\) in a different such model shows
that existing arguments need not give an existing function value.

Under an inner-valued free assignment \(g(x)\in D\),
\(E!(x)\) is true in every model.  The outer assignment
\(g(x)=1\) in the same permitted model falsifies it.
This proves that an unspecified free-variable range cannot
determine which of these open-formula consequence relations is meant.

No existing object satisfies falsum, so the description
\(\iota x\bot\) has no value.  By the explicit negative
description policy an atom with this complete argument, including its
existence and self-identity atoms, is false; ordinary function applications
with this undefined input remain undefined and yield the same atomic
failure. Classical negation of such an atom is true. This does not assert
failure of an outer description whose formula body merely mentions that
term, since that formula body can evaluate a negated false atom as true.
For failed uniqueness choose a model with \(D=U=\{0,1\}\).
Both existing objects satisfy \(E!(x)\), so
\(\iota xE!(x)\) has no unique value and is undefined.
The same atomic clause then applies.  These are genuinely
undefined terms, unlike the defined outer constant of the first model.

If one changed the convention to \(D=\varnothing\) while
retaining nonempty \(U\), no witness could satisfy
\(\exists xE!(x)\).  Moreover the premise
\(\{E!(y)\}\models A\) would hold vacuously for any \(A\),
including falsum; the nonempty-domain row would then incorrectly
produce an unconditional consequence of falsum.  This verifies
both failures and the necessity of excluding that different model class.
\end{proof}
\end{annexitem}

\begin{annexitem}{annex:p8.negative.certificate-example}{thm:p08-negative-certificate-example}{A complete finite target proof with value witnesses}
\begin{proof}
Use the translation clauses with bound auxiliaries consistently renamed.
The two target premises are
\[
 H=\forall x\left(E(x)\to
      \exists v(v=x\land E(v)\land P(v))\right),\qquad
 J=\exists u(u=c\land E(u)).
\]
The target conclusion is
\(K=\exists v(v=c\land E(v)\land P(v))\).
These are exactly the translations of the displayed source premises and
conclusion, not additional target axioms. Choose target variables \(a,b\)
fresh for \(H,J,K\) and for each other's witness subproof.

Eliminate \(J\) by temporarily assuming
\(a=c\land E(a)\). Conjunction elimination gives \(a=c\) and
\(E(a)\). Classical universal elimination in the target, at the target
variable \(a\), gives from \(H\)
\[
 E(a)\to\exists v(v=a\land E(v)\land P(v)).
\]
Implication elimination with \(E(a)\) gives its existential consequent.
Eliminate this existential by temporarily assuming
\(b=a\land E(b)\land P(b)\), extracting \(b=a,E(b),P(b)\) by
conjunction elimination. From \(a=c\) and \(b=a\), target equality
elimination, applied to the formula \(b=x\), gives \(b=c\).
Conjunction introduction now gives \(b=c\land E(b)\land P(b)\),
with its parentheses matching the translation's chosen association.
Classical existential introduction at target term \(b\) yields \(K\).

The inner existential elimination discharges the complete conjunction
assumption for \(b\). Its eigenvariable occurs in neither \(K\) nor
the remaining premises \(H,J,a=c\land E(a)\), so the rule's condition
holds. The outer existential elimination then discharges the complete
conjunction assumption for \(a\); its eigenvariable occurs in neither
\(K\) nor \(H,J\). Both conjunction assumptions were actually used,
so these eliminations also satisfy the registered nonvacuous discharge
discipline. Copies of \(H\) or an extracted conjunct use derivation
label names renamed fresh when required. This is a finite proof on
Sheet~\ref{sheet:nd-fol} from \(H,J\) of \(K\). Its assumptions are
therefore drawn from \(\Theta\cup\delta[\{\forall xP(x),E!(c)\}]\),
and the certificate row proves the claimed source judgment. The closed
nonemptiness axiom is permitted but not needed in this particular proof.
\end{proof}
\end{annexitem}
